\chapter{Emergence and novelty}

\section{Two words that need disentangling}

The Game of Life produces gliders. The flock of Boids has a shape. The
Tierra ecology grows parasites. Each of these is sometimes called
\emph{emergent}, and each is sometimes called \emph{novel}. The two words
are not synonyms, and the literature does not always agree on the
difference. This chapter is about what they each mean, where they are
made precise, and which of the precise versions the lens will commit to.

We will not pretend the disagreements collapse. Emergence and novelty are
both contested terms, with several technically respectable definitions
each, and parts of the field are arguing about which one is right. The
book takes a clear position --- Hoel-style causal emergence for emergence,
observer-relative divergence for novelty --- and explains the alternatives
in enough detail that you can read papers that adopt them and know what
you are reading.

\section{Weak emergence: in-principle-derivable, in-practice-surprising}

The mildest version of emergence is what Mark Bedau called \emph{weak
emergence}. A property of a system is weakly emergent if, given the
rules, you cannot in practice predict the property from the rules without
running the system. In principle you could; in practice you could not.
The gliders of the Game of Life are weakly emergent in this sense: the
rule for cell updates does not mention gliders, and there is no shortcut
that lets you derive ``gliders exist'' from the rule, even though
gliders exist whenever you let the rule run on the right initial
conditions.

This is a defensible definition. It pins down a real phenomenon ---
patterns that are computationally irreducible to their generating rule
--- and it does so without committing to any metaphysics. Weak emergence
does not assert that the gliders are anything more than gliders. It
asserts only that they are not derivable in advance.

The two reservations: weak emergence does not, by itself, give you a
quantitative criterion (how weakly emergent is something?), and it
inherits a notion of ``in-practice unpredictable'' that depends on the
computational resources of the predictor. A pattern that is weakly
emergent for a person with pencil and paper may be entirely transparent
to a supercomputer. The intuition is real but the formalism is
soft.

\section{Strong emergence: the controversial one}

\emph{Strong emergence} asserts more: that a property of a system is not
even \emph{in principle} derivable from the properties of the parts.
Strong emergence has been claimed for consciousness, for life, and for a
handful of other big-ticket properties. It is philosophically contentious
and largely outside the scope of artificial life --- the systems we build
are computational, and their behaviour is derivable from the rules even
if computing the derivation is expensive.

For the book, strong emergence is a vocabulary item rather than a working
definition. We will mention it twice more: once in the context of
philosophical questions about whether artificial-life creatures can be
``really alive,'' and once in the context of arguments that some systems
exhibit properties their substrate alone cannot account for. Neither
discussion will hinge on whether strong emergence exists.

\section{Supervenience: the boring but useful relation}

Supervenience is the formal companion to talk about emergence. A property
$A$ supervenes on a property $B$ if there cannot be a difference in $A$
without a difference in $B$. The shape of a flock supervenes on the
positions of its members: there cannot be two flocks with the same set
of positions but different shapes. The gliders of the Game of Life
supervene on the cell states: there cannot be two GoL grids with the
same cell states but different glider populations.

Supervenience is too weak to be a definition of emergence (anything
supervenes on its parts), but it is strong enough to be useful: it lets
us talk about a higher-scale property without committing to its existence
as anything beyond a re-description of the lower scale. The lens uses
supervenience implicitly: the scale projection $\pi$ takes a
configuration of substrate sites and produces a feature (centroid, mass,
descriptor). The feature supervenes on the sites. That is what makes the
feature a feature \emph{of the configuration} and not an independent
fact.

\section{Causal emergence: the version the book will use}

The definition of emergence the book commits to is due to Erik Hoel and
collaborators, and it is built on a quantity called \emph{effective
information}. The idea is unusual enough to deserve a careful
walk-through.

\subsection*{Effective information, in three sentences}

Take a system that is a finite-state Markov chain. Imagine intervening on
its current state uniformly at random --- equally likely to start the
system in any state --- and ask how much information that intervention
gives you about its next state. The amount of information you get is the
effective information.

That is the definition. The rest is unpacking what the words mean. Two
things matter:

\begin{itemize}[leftmargin=1.5em,itemsep=2pt]
  \item \textbf{Determinism.} If the transition rule is very deterministic
    --- the next state is nearly a function of the current state --- the
    intervention is informative, because intervening tells you almost
    exactly what comes next.
  \item \textbf{Non-degeneracy.} If different starting states tend to
    flow to different next states (rather than all collapsing to the
    same attractor), the intervention is also informative, because
    different interventions give you different futures.
\end{itemize}

A maximally effective-informational system is both deterministic and
non-degenerate: each input maps cleanly to its own output. A maximally
ineffective system is either pure noise (next state independent of
current) or maximally collapsing (all roads lead to the same place). The
exact formula appears in the math primer at the back of the book; for
now, take effective information as a number that measures how
informative interventions are.

\subsection*{Causal emergence as a difference}

Now imagine you have a coarse-graining of the system: a way of grouping
states into macro-states. The Game of Life has many candidate
coarse-grainings: ``how many live cells are there'', ``where is the
centroid of the live region'', ``is there a glider present.'' Each
coarse-graining gives you a new Markov chain --- the macro chain --- by
averaging the micro-transitions inside each macro-state.

The Hoel definition: \emph{a system is causally emergent at a
coarse-graining $\Pi$ if the macro chain has higher effective
information than the micro chain.} That is the whole of it. Causal
emergence is the proposition that some choice of macroscale carries more
predictive information per intervention than the microscale you started
with.

The intuition is sharper than the formula. The Game of Life at the cell
level is fairly chaotic; many initial states lead to similar
configurations after enough steps. But if you look at it at the
glider-level --- ``how many gliders are travelling north-east, with what
phase'' --- the dynamics are very deterministic. Effective information
\emph{goes up} when you switch to the glider description. That increase
is what Hoel and collaborators call causal emergence.

\subsection*{Why this definition rather than another}

Three properties:

\begin{itemize}[leftmargin=1.5em,itemsep=2pt]
  \item \textbf{Quantitative.} You get a number, not a yes/no.
  \item \textbf{Intrinsic.} The number is computed from the transition
    matrix; no observer's preferences enter.
  \item \textbf{Coarse-graining-explicit.} The definition forces you to
    name the macroscale you are testing. Two researchers who disagree
    about whether a system is emergent can be made to disagree about
    \emph{which coarse-graining}, which is a more productive
    disagreement.
\end{itemize}

The book takes effective information as the default definition of
emergence. We will not pretend the choice is uncontroversial. The
intervention distribution (uniform over states) is a convention rather
than a derivation, and results are sensitive to it. The supremum-over-
coarse-grainings is uncomputable for large state spaces. $\Phi$-ID --- an
alternative framework from Rosas, Mediano, and collaborators --- replaces
partitions with arbitrary supervenient features, and asks whether the
macro-feature has \emph{unique} predictive information about the future
that no microscopic component has. $\Phi$-ID is feature-agnostic and
methodologically more demanding; the book does not develop it in detail.

For working purposes, treat causal emergence as the technical version
of ``a higher-scale description is more predictive than the lower-scale
description it summarises.'' That is what the gliders give you in the
Game of Life. It is also, plausibly, what the flock gives you in Boids
and what the parasite ecology gives you in Tierra. Whether these
intuitions can be made formal is, in each case, open: see OP1 in
Chapter~9.

\section{Computational irreducibility: a different vocabulary for the same place}

Stephen Wolfram's contribution to the discussion is the phrase
\emph{computational irreducibility}: some systems' future behaviour
cannot be computed without simulating them step by step. The connection
to weak emergence is direct. A weakly emergent property is one that
requires you to run the simulation. Computational irreducibility is the
property of the system that forces you to.

Wolfram's classification of cellular automata into four classes ---
quiescent, periodic, chaotic, complex --- is a regime-of-behaviour
taxonomy. Class IV (complex) is the ``edge of chaos'' regime where
non-trivial structures persist and interact. The Game of Life is Class
IV. Lenia parameter sets that produce creatures are Class IV.

The lens does not adopt Wolfram's vocabulary as primary, but the
classification is useful for one purpose: it points at the kinds of
substrates where the book's emergence question is most interesting. A
Class I CA (quiescent) has no emergence in any sense. A Class III CA
(chaotic) has too much. The cases worth studying live in the middle.

\section{Novelty, novelty search, and the trouble with the word}

Switch tracks. \emph{Novelty} is the second contested word. In the
loosest sense, novelty is the production of something that did not
exist before. The Game of Life produces novelty trivially --- every new
configuration is one that did not exist last tick. Nobody calls that
novelty in the artificial-life sense. The interesting kind of novelty is
\emph{something that did not exist before and is qualitatively
different}, and the qualitatively-different part is where the
definitions diverge.

\subsection*{Novelty search}

Joel Lehman and Kenneth Stanley's \emph{novelty search} made the slogan
precise inside an evolutionary algorithm. The fitness function rewards
behavioural distance from an archive of previously-seen behaviours
rather than progress toward an objective. The archive grows as new
behaviours appear. The result, in their experiments, was that novelty
search could outperform objective-driven search on hard deceptive
problems --- precisely because following the objective tends to get
stuck on local optima, and a novelty-driven walk does not.

The technical content of novelty search is a divergence measure: how far
is a new behaviour from its $k$ nearest neighbours in the archive?
Behaviours with high divergence get to survive and reproduce; behaviours
that crowd into an already-populated part of behaviour space get culled.
Novelty is observer-relative in the strict sense: it depends on the
descriptor space the observer chose, and on the archive.

Recent analysis has clarified that novelty search is best read at the
level of archive coverage and saturation. The slogan ``objectiveless
evolution'' overstates the case --- novelty search \emph{does} have an
objective, in the form of archive coverage --- but the underlying point
survives: rewarding diversity rather than improvement can find things
that improvement-following misses.

\subsection*{Quality-Diversity}

Novelty search was the start of a family. The Quality-Diversity (QD)
algorithms --- MAP-Elites being the canonical instance --- maintain a
repertoire of solutions that are simultaneously diverse and
high-performing on some objective. Recent QD work on Lenia explicitly
uses this approach to discover diverse self-organising patterns. The
observer keeps a behaviour grid and stores the best-performing genome in
each cell.

QD's relationship to novelty search is straightforward: novelty search
optimises only for coverage, QD optimises for coverage \emph{and}
quality. In the lens, both are observer-state choices --- archive,
descriptor map, update rule.

\subsection*{Open-endedness}

\emph{Open-endedness} is the larger question that novelty search
gestures at: not just \emph{will the system produce new things}, but
\emph{will it keep doing so, indefinitely}. This is the property natural
evolution seems to have and that artificial systems mostly do not.
Tierra produced novelty for hours and then stopped. POET co-evolves
environments and agents to push them both toward harder problems and
has produced longer runs of novelty, but ``indefinite'' remains an
aspiration rather than a demonstrated fact.

Several technical definitions of open-endedness compete:

\begin{itemize}[leftmargin=1.5em,itemsep=2pt]
  \item \textbf{Hughes--Dennis learnable novelty} --- the rate of
    \emph{learnable} new skills is bounded away from zero in the long
    run.
  \item \textbf{The $\Omega$ metric} --- a residence-time--based measure
    of how long the system stays in any given region of descriptor
    space. Open-endedness shows up as long-tailed residence
    distributions.
  \item \textbf{Tria-style adjacent-possible} --- the rate of new-class
    arrivals follows Heaps' law $\sim t^\beta$ with $\beta<1$, generated
    by P\'olya-urn-with-triggering dynamics where finding new things
    makes more new things findable.
\end{itemize}

These are not the same. Each picks a different aspect of
``indefinitely-novel.'' The book does not pick one. The observer slot is
where they live, and the system description should record which observer
was used.

\subsection*{Bounded versus unbounded novelty}

A useful distinction: a system can produce novelty without producing
unbounded novelty. Hazen and Wong's work on mineral evolution is an
example. Minerals on Earth evolve --- the set of mineral species grows
over geological time, functional information increases --- but the
process is bounded by the chemistry of the planet. There are only so
many ways to combine the elements. The bounded case is interesting in
its own right: it shows what domain-bounded novelty looks like, and it
sharpens the question of whether unbounded novelty (genuine
open-endedness) is a real possibility or an aspirational ideal.

The book's stance: the bounded case is where most artificial-life
systems live, and is worth taking seriously. Unbounded novelty is an
open frontier (Chapter~9).

\section{Putting the words to work}

We have surveyed two clouds of meaning. Emergence: from supervenience
through weak emergence through computational irreducibility to causal
emergence to $\Phi$-ID. Novelty: from raw difference through novelty
search through Quality-Diversity to open-endedness in its several
technical senses. The vocabulary is large and not entirely consistent.

The book's commitments, restated:

\begin{itemize}[leftmargin=1.5em,itemsep=2pt]
  \item \textbf{Emergence} = Hoel-style causal emergence: a positive
    difference in effective information between a macroscale and a
    microscale. The book will sometimes use the word weakly to mean
    ``something macroscopic and unexpected happened,'' but when we want
    to be precise we mean causal emergence with a stated coarse-graining.
  \item \textbf{Novelty} = observer-relative divergence: the score
    assigned by an observer to a part of the trajectory that has not
    been seen before in the observer's representation.
  \item \textbf{Open-endedness} = a property of trajectories under a
    specified observer: novelty (in the divergence sense) is not
    eventually saturated. The book treats this as observer-relative by
    default; the question of substrate-intrinsic open-endedness is
    open.
\end{itemize}

These commitments slot into the five-slot lens cleanly. Emergence is
about $M(\mathcal{S})$, the Markov chain induced by the joint update
rule. Novelty is about $O$, the observer. Open-endedness is about the
long-time behaviour of $O$. Chapter~4 builds the lens; Chapter~8 returns
to emergence and novelty quantitatively, with the formulas spelled out.

\section*{To think about}

\begin{enumerate}
  \item Find a system you would describe as ``emergent.'' Write down two
    distinct candidate coarse-grainings of it. Which one would you bet
    has the higher effective information? Why?
  \item Novelty search is sometimes described as ``objectiveless.'' The
    book has argued it is not. What \emph{is} the objective?
  \item Tierra ran for several hours and then its novelty rate dropped
    to near zero. Pick one of the three open-endedness frames in this
    chapter and describe what changed in Tierra at the point of
    saturation.
  \item Strong emergence is contested in philosophy. Suppose you were
    asked to argue that a Lenia creature is strongly emergent. What
    would your strongest argument be? What would the most natural
    rebuttal be?
  \item Computational irreducibility is a property of the rule; causal
    emergence is a property of the system at a coarse-graining. Are
    they related, in the sense that one implies the other? If yes, in
    which direction?
\end{enumerate}
